\section{Architecture and Implementation}
\label{sec:implementation}
Afar-prgenv integrates AMD compiler drops into the Cray PE by providing a
module hierarchy and lightweight wrapper logic. The overlay is intentionally
minimal: it avoids replacing the PE and instead augments it with AFAR-specific
paths, metadata, and diagnostics.

\begin{figure*}[t]
  \centering
  \begin{tikzpicture}[
  box/.style={draw, rounded corners, align=center, minimum width=2.6cm, minimum height=0.9cm},
  small/.style={draw, rounded corners, align=center, minimum width=2.2cm, minimum height=0.8cm},
  arrow/.style={-Latex, thick}
]
  \node[box] (user) {User build\\(CMake/Make)};
  \node[box, right=1.8cm of user] (wrappers) {Cray wrappers\\\texttt{ftn/cc/CC}};
  \node[box, right=1.8cm of wrappers] (compiler) {AFAR compiler\\drop};

  \node[small, below=1.2cm of wrappers] (pkg) {pkg-config\\metadata};
  \node[small, below=1.2cm of compiler] (rocm) {ROCm toolkit};

  \node[box, below=1.2cm of user] (modules) {Afar-prgenv\\modules};

  \draw[arrow] (user) -- (wrappers);
  \draw[arrow] (wrappers) -- (compiler);
  \draw[arrow] (wrappers) -- (pkg);
  \draw[arrow] (compiler) -- (rocm);
  \draw[arrow] (modules) -- (wrappers);
  \draw[arrow] (modules) -- (compiler);
  \draw[arrow] (modules) -- (rocm);
\end{tikzpicture}

  \caption{Afar-prgenv overlay architecture. The Cray PE remains the primary
  interface; Afar-prgenv inserts AFAR compilers and ROCm while preserving PE
  wrapper semantics and library metadata.}
  \Description{Block diagram showing user builds flowing through Cray wrappers
  into AFAR compilers, with Afar-prgenv modules providing ROCm and pkg-config
  metadata.}
  \label{fig:architecture}
\end{figure*}

\subsection{Module Hierarchy}
\label{sec:modules}
The module tree is split into a meta module (\texttt{afar-prgenv/<ver>}) and two
component modules: \texttt{afar-rocm/\allowbreak <ver>} and
\texttt{afar-amd/\allowbreak <ver>}. The meta module
unloads conflicting vendor modules, then loads the ROCm toolkit module followed
by the compiler module to ensure the compilers are last in `PATH`. This keeps
the Cray PE wrappers stable while redirecting the underlying toolchain.

\subsection{Compiler Wrappers}
\label{sec:wrappers}
Afar-prgenv ships `ftn`, `cc`, and `CC` wrappers that preserve Cray PE behavior
while injecting AFAR-specific fixes. The wrappers:
\begin{itemize}
  \item Resolve the real Cray compiler driver and delegate to it.
  \item Clear \texttt{CRAY\_ACCEL\_*} to avoid unsupported PE offload flags.
  \item Inject `--offload-arch` when `-fopenmp` is present and no offload arch
  is explicitly specified.
  \item Re-synchronize MPI metadata at invocation time to match the active
  `cray-mpich` module.
  \item Filter or replace incompatible flags using a configurable rule set.
\end{itemize}

\begin{table}[t]
  \centering
  \small
  \setlength{\tabcolsep}{4pt}
  \caption{Selected Afar-prgenv environment overrides.}
  \label{tab:env-overrides}
  \begin{tabular}{l l}
    \hline
    Variable & Purpose \\
    \hline
    \texttt{AFAR\_REAL\_FTN} & Override real PE `ftn` path \\
    \texttt{AFAR\_REAL\_CC} & Override real PE `cc` path \\
    \texttt{AFAR\_REAL\_CXX} & Override real PE `CC` path \\
    \texttt{AFAR\_FTN\_OFFLOAD\_ARCH} & Force Fortran offload arch \\
    \texttt{AFAR\_WRAPPER\_FILTER} & Custom wrapper filter file \\
    \texttt{AFAR\_WRAPPER\_FILTER\_EXTRA} & Add local filter rules \\
    \texttt{AFAR\_PKGCONFIG\_}\allowbreak\texttt{SHIM\_DIR} & Extra \texttt{.pc} shim path \\
    \texttt{AFAR\_PKGCONFIG\_}\allowbreak\texttt{SHIM\_AUTO} & Auto-refresh shims toggle \\
    \texttt{AFAR\_LLVM\_LIB\_DIR} & Override LLVM runtime fallback \\
    \hline
  \end{tabular}
\end{table}

\subsection{Wrapper Filtering and Link Type Guard}
\label{sec:wrapper-filter}
Wrapper filters ignore or replace arguments injected by build systems that are
incompatible with AFAR. The filter is defined in a shared rule file and can be
extended per site via \texttt{AFAR\_WRAPPER\_}\allowbreak\texttt{FILTER\_EXTRA}. The wrappers also
guard against Cray's dynamic link mode by unsetting
\texttt{CRAYPE\_LINK\_TYPE=dynamic}, preventing late-injected flags that
conflict with the AFAR toolchain.

\begin{table}[t]
  \centering
  \small
  \setlength{\tabcolsep}{4pt}
  \caption{Default wrapper filter rules used by Afar-prgenv.}
  \label{tab:wrapper-filter}
  \begin{tabular}{l p{0.6\columnwidth}}
    \hline
    Pattern & Rationale \\
    \hline
    \texttt{-dynamic} & Avoid PE-injected dynamic link flags that break AFAR. \\
    \texttt{-ftree-vectorize} & Drop GCC-specific flag not supported by AFAR. \\
    \texttt{-ffunction-sections} & Drop GCC-specific flag in PE build flags. \\
    \texttt{-fdata-sections} & Drop GCC-specific flag in PE build flags. \\
    \hline
  \end{tabular}
\end{table}

\subsection{pkg-config Integration}
\label{sec:pkgconfig}
MPI and vendor libraries depend on consistent `pkg-config` metadata. The
generator produces AFAR-specific `.pc` files for MPI and ROCm, while a shim
generator captures module-provided metadata for Cray libraries that do not ship
usable `.pc` files. The wrappers mirror the PE-specific metadata path into
\texttt{PKG\_CONFIG\_PATH} to keep build systems stable under module swaps.
When required ROCm 6 runtime libraries are missing, the AFAR ROCm module
suppresses Cray MPICH GTL `pkg-config` variables to avoid injecting
incompatible link flags. This keeps MPI builds functional while signaling that
accelerator-enabled MPI paths need follow-up validation.

\begin{figure*}[t]
  \centering
  \resizebox{\textwidth}{!}{\begin{tikzpicture}[
    box/.style={draw, rounded corners=2pt, align=center, text width=2.4cm, minimum height=0.85cm},
    arrow/.style={-Latex, thick},
    font=\small
  ]
  \node[box] (pe) {Cray PE modules\\(cpe/PrgEnv/MPICH)};
  \node[box, right=1.1cm of pe] (fixed) {PE\_AMD\_FIXED\\\_PKGCONFIG\_PATH};
  \node[box, right=1.1cm of fixed] (wrap) {Afar-prgenv\\wrappers};
  \node[box, right=1.1cm of wrap] (pkg) {PKG\_CONFIG\_PATH};
  \node[box, right=1.1cm of pkg] (build) {pkg-config\\and builds};

  \node[box, below=1.2cm of wrap] (shim) {Shim \texttt{.pc} files\\(\texttt{AFAR\_PKGCONFIG\_}\allowbreak\texttt{SHIM\_DIR})};

  \draw[arrow] (pe) -- (fixed);
  \draw[arrow] (fixed) -- (wrap);
  \draw[arrow] (wrap) -- (pkg);
  \draw[arrow] (pkg) -- (build);
  \draw[arrow] (shim) -- (wrap);
\end{tikzpicture}
}
  \caption{pkg-config metadata flow under Afar-prgenv.}
  \Description{Diagram showing PE metadata paths mirrored into PKG_CONFIG_PATH,
  with AFAR shim directories appended by the wrappers.}
  \label{fig:pkgconfig-flow}
\end{figure*}

Shims can be regenerated automatically when the module environment changes, so
libraries can be loaded before or after
\texttt{afar-}\allowbreak\texttt{prgenv} without breaking builds. The wrappers
track environment changes and re-run shim generation when necessary, keeping
\texttt{PKG\_}\allowbreak\texttt{CONFIG\_PATH} consistent for build systems.

\subsection{MPICH Flavor Mapping}
\label{sec:mpich-flavor}
Cray MPICH 8.x uses \texttt{mpich3.4a2} and 9.x uses \texttt{mpich4.3.1}.
Afar-prgenv generates per-flavor `mpichf90.pc` files and selects the correct
flavor based on
\texttt{CRAY\_MPICH\_VERSION}. The wrappers re-check this at compile time so MPI
metadata remains correct even if users swap MPICH modules after loading AFAR.

\subsection{LLVM Runtime Fallback}
\label{sec:llvm-runtime}
Some AFAR drops omit LLVM runtime libraries. Examples include
\texttt{libpgmath.so} and \texttt{libflang.so}. The module logic detects the newest compatible
\texttt{/opt/rocm-*/llvm/lib} and appends it as a fallback, unless a site
overrides the path with \texttt{AFAR\_LLVM\_LIB\_DIR}. This preserves compiler
functionality without modifying the upstream drop.

\subsection{Generation and Configuration}
\label{sec:generator}
Version mappings and PE paths are defined in configuration files and are
regenerated with a single script. This keeps the module tree reproducible and
aligns the AFAR drop version with the ROCm toolkit and MPI metadata used in
tests. The generator also writes diagnostics that the support team uses to
track drop readiness over time.
